\section{Исходный код}
Реализация состоит из двух частей: \textbf{сбор данных и веб-интерфейс (Python)} и \textbf{обработка/поиск (C++)}. Все компоненты разработаны с акцентом на производительность, простоту и минимальные зависимости. Система не использует STL за пределами \texttt{std::string} и \texttt{std::vector}, а также сторонние поисковые библиотеки.

\subsection*{Сбор данных (\texttt{crawler.py})}
Модуль загружает архивы \texttt{.7z} с сайта Archive.org, содержащие дампы сообществ Stack Exchange:
\begin{itemize}
\item \texttt{ai.stackexchange.com}
\item \texttt{datascience.stackexchange.com}
\end{itemize}
После скачивания архив распаковывается с помощью библиотеки \texttt{py7zr}, извлекается файл \texttt{Posts.xml}. Затем парсер \texttt{xml.etree.ElementTree} проходит по всем записям типа \texttt{PostTypeId="1"} (вопросы), извлекая заголовок и тело сообщения.

HTML-разметка очищается с помощью регулярных выражений:
\[
\texttt{re.sub(r"<[^>]+>", "", html)}
\]
Текст объединяется: \texttt{title + "\textbackslash n\textbackslash n" + body}, и проверяется на минимальную длину — не менее 100 слов. Каждый документ сохраняется в директорию \texttt{corpus/} как \texttt{<source>\_<id>.txt}, а метаданные — в JSON-файл с тем же именем, включая:
\begin{itemize}
\item URL статьи;
\item Хеш содержимого (MD5);
\item Время извлечения;
\item Количество слов.
\end{itemize}
Для обеспечения идемпотентности повторная загрузка не перезаписывает существующие файлы. Между запросами соблюдается задержка 1 секунда.

\subsection*{Токенизатор (\texttt{tokenizer.cpp})}
Работает в два этапа:
\begin{enumerate}
\item Чтение всего файла в память через \texttt{std::ifstream} в двоичном режиме.
\item Посимвольная обработка:
\begin{itemize}
\item Латинские буквы приводятся к нижнему регистру;
\item Допустимыми символами считаются только \texttt{[a–z0–9]};
\item Прочие символы являются разделителями.
\end{itemize}
\end{enumerate}
Накопленное слово добавляется в список, если его длина ≥2. Результат — вектор нормализованных терминов, готовых к стеммингу и индексации.

\subsection*{Стеммер (\texttt{PorterStemmer.h})}
Реализован упрощённый алгоритм Портера на C++. Основные шаги:
\begin{itemize}
\item Удаление суффиксов: \texttt{s}, \texttt{ss}, \texttt{ies}, \texttt{ed}, \texttt{ing};
\item Правила замены: \texttt{ational → ate}, \texttt{tional → tion}, \texttt{ization → ize};
\item Удаление окончаний: \texttt{al}, \texttt{ance}, \texttt{ence}, \texttt{er}, \texttt{ic};
\item Коррекция: \texttt{e} удаляется при определённых условиях;
\item Удвоение согласных: например, \texttt{ll} после \texttt{e} может быть сокращено до \texttt{l}.
\end{itemize}
Функция \texttt{stem()} принимает строку, приводит её к нижнему регистру и последовательно применяет правила. Алгоритм работает за $O(k)$ для слова длины $k$.

\subsection*{Построение индексов (\texttt{index.cpp})}
Обрабатываются все файлы в директории \texttt{corpus/}. Для каждого документа:
\begin{itemize}
\item Извлекаются токены;
\item Каждый токен стеммируется;
\item Формируется обратный индекс: \texttt{term → [doc\_id]}.
\end{itemize}
Также строится прямой индекс: \texttt{doc\_id → title, url}, где URL формируется автоматически на основе имени файла.

Индексы сериализуются в бинарные файлы:
\begin{itemize}
\item \texttt{inverted\_index.bin}: сначала количество терминов, затем для каждого:
\begin{itemize}
\item Длина строки (4 байта);
\item Само слово;
\item Количество документов (4 байта);
\item Список ID документов (по 4 байта каждый).
\end{itemize}
\item \texttt{forward\_index.bin}: аналогично — число документов, затем для каждого:
\begin{itemize}
\item \texttt{doc\_id}, длина и текст заголовка, длина и URL.
\end{itemize}
\end{itemize}
Перед сохранением обратный индекс сортируется по алфавиту. Это позволяет использовать бинарный поиск при выполнении запросов.

\subsection*{Булев поиск (\texttt{search.cpp})}
Поддерживается полный синтаксис:
\begin{itemize}
\item \texttt{car engine} → \texttt{car \&\& engine}
\item \texttt{car || truck}
\item \texttt{!bike}
\item \texttt{(car || truck) \&\& engine}
\end{itemize}
Архитектура:
\begin{enumerate}
\item Загрузка индексов из бинарных файлов;
\item Лексический анализ: входной запрос разбивается на токены (слова, \texttt{&&}, \texttt{||}, \texttt{!}, скобки);
\item Синтаксический анализ методом рекурсивного спуска:
\begin{itemize}
\item \texttt{eval\_expr()} — обрабатывает OR;
\item \texttt{eval\_term()} — AND;
\item \texttt{eval\_factor()} — NOT и скобки.
\end{itemize}
\item Операции над списками:
\begin{itemize}
\item \texttt{AND}: слияние двух отсортированных списков за $O(|A|+|B|)$;
\item \texttt{OR}: объединение с удалением дубликатов;
\item \texttt{NOT}: дополнение до полного множества документов.
\end{itemize}
\end{enumerate}
Поиск выполняется за доли миллисекунды даже на корпусе из 30\,000 документов.

\subsection*{Веб-интерфейс (\texttt{web.py})}
Для удобства используется веб-интерфейс на Flask. Он:
\begin{itemize}
\item Отображает форму ввода запроса;
\item При поиске запускает внешнюю программу \texttt{./searcher "запрос"} через \texttt{subprocess.run()};
\item Парсит вывод: каждая строка вида \texttt{title | url} преобразуется в ссылку;
\item Поддерживает постраничную навигацию (по 50 результатов на страницу).
\end{itemize}
Интерфейс работает на \texttt{http://localhost:5000}, имеет чистый, отзывчивый дизайн и полностью автономен.

Производительность:
\begin{itemize}
\item Время ответа: $\sim$1–5 мс;
\item Поддержка сложных логических выражений;
\item Масштабируемость: легко адаптируется под новые источники.
\end{itemize}

\pagebreak
